\subsubsection{二回分の更新}\label{節：周期係数を行列で見直す}
  \sectionref*{節：周期的な係数}では,$2x+1$の後に$3x+2$を計算し,$6x+5$を得た.
  今度は,それぞれの更新を行列で書いてみよう.
  \[
    A=\begin{pmatrix}2&1\\0&1\end{pmatrix},\qquad
    B=\begin{pmatrix}3&2\\0&1\end{pmatrix}.
  \]
  先に$A$,次に$B$で更新するので,二回分は$BA$になる.
  実際に掛けると,
  \[
    BA=\begin{pmatrix}6&5\\0&1\end{pmatrix},\qquad
    BA\binom{x}{1}=\binom{6x+5}{1}
  \]
  となり,本文の計算に戻る.順序を逆にすると,
  \[
    AB=\begin{pmatrix}6&5\\0&1\end{pmatrix}.
  \]
  この例ではたまたま同じ行列になった.一般にも同じとは限らない.
  接続演習問4の更新を使うと,
  \[
    U=\begin{pmatrix}-1&0\\0&1\end{pmatrix},\qquad
    V=\begin{pmatrix}2&1\\0&1\end{pmatrix}
  \]
  に対して,
  \[
    VU=\begin{pmatrix}-2&1\\0&1\end{pmatrix},\qquad
    UV=\begin{pmatrix}-2&-1\\0&1\end{pmatrix}
  \]
  となる.更新を行う順序が,付加項の符号に現れた.
  一般の$a_{n+1}=p_na_n+q_n$でも,$\binom{a_{n+1}}1=
  \begin{pmatrix}p_n&q_n\\0&1\end{pmatrix}\binom{a_n}1$と表せる.
  二回の更新は,先に行う\bookterm{matrix}{行列}を右に置いた積になる.
  一周期分をまとめる操作は,その周期の行列をこの順序で掛けることに当たる.
  途中の値は,本文と同じように一回ずつの更新から求められる.
  \bookproblemref{practice-connections}{4}{接続演習問4}の式と比べると,
  どちらの表現でも同じ更新を追っていることが分かる.
